\subsection{The Shared Interaction Layer}\label{subsec:interaction_layer}

Every scene is built on the same foundation. A scene component reads the
receiver's skeleton and gestures each frame and acts on them, and the gestures
themselves are the discrete labels emitted by MediaPipe's
recogniser~(\autoref{subsec:mediapipe}), of which the application uses five ---
\texttt{Closed\_Fist}, \texttt{Pointing\_Up}, \texttt{Thumb\_Up},
\texttt{Thumb\_Down}, and \texttt{Victory}. These form a small command set, reused
consistently across navigation, menus, and direct manipulation, so that a gesture
learned in one scene carries a familiar role in the next.

Movement between scenes is itself gestural. \texttt{CircularSceneChanger} cycles
the four scenes in a fixed order, and the two thumb gestures held in opposition
advance or rewind the cycle: left \texttt{Thumb\_Down} with right \texttt{Thumb\_Up}
moves forward, and the reverse combination moves back. A short cooldown guards
against repeated triggering, and because the timer is held in a static field it
survives the scene reload, preventing the gesture that loaded a scene from
immediately loading the next. The application is therefore navigated end to end
without a pointer or a menu button.

A second mechanism recurs wherever a gesture must be held over time. The detector
occasionally drops a gesture for a frame or two, re-triggering the palm stage when
its confidence falls~(\autoref{subsec:mediapipe}); taken literally, a held grip
would flicker. Each grasp-based interaction therefore tolerates a brief lapse:
once a \texttt{Closed\_Fist} is lost, a short release timer must elapse before the
grip is actually let go, so a momentary dropout is treated as continued holding.
The same short-delay pattern appears in the face deformer, the grappling bear, and
the omni-directional movement controller, and is described once here rather than
repeated in each.

Scenes that offer several modes select between them through a body-anchored menu
rather than a screen overlay. \texttt{BraceletMenuCore} renders a wristband that
snaps to the right wrist each frame, encircled by three option spheres. A
right-hand \texttt{Victory} toggles the menu into view; the user then rotates the
wrist until the desired option rises to the top, where the highest sphere drops
beneath the band, enlarges, and shows a floating label. Selection is the left
hand's \texttt{Thumb\_Up}, and both hands' \texttt{Thumb\_Down} resets to the
default. Colour marks state --- one tint for an idle option, another for the one
currently hovered, a third for the mode actually running. The menu is diegetic: it
lives on the hand, requires no controller and no two-dimensional interface, and
reuses the same gestures as everything else. The base class provides this
behaviour and leaves the binding of options to actions to a per-scene subclass, so
that the city and painting scenes share one menu mechanism with different contents.

\begin{figure}[H]
	\centering
	\captionsetup{justification=centering}
	\includegraphics[width=0.9\linewidth, frame]{img/bracelet_menu.png}
	\caption{The wrist bracelet menu, with the hovered option
	         raised below the band and labelled.}
	\label{fig:bracelet-menu}
\end{figure}

Pointing at the world is also shared. \texttt{HandPointer} casts a ray and reports
what it strikes, in either of two aiming styles chosen by a right-hand gesture. In
the pointing style the ray leaves the index fingertip along the finger's direction;
in the palm style it leaves the centre of the palm along the palm's normal, the
latter computed as the cross product of the vectors from the wrist to the index and
little-finger knuckles. A beam and a marker show the aim, and a left-hand
\texttt{Victory} signals confirmation. The companion bear and the city's building
editor both build on this one component, reading its hit point and its confirmation
state rather than deriving a ray of their own.
